\section{Discussion}
Structural metrics are the primary empirical success criterion because they measure whether usable-update mass matches the public deployment target. Improved alignment need not imply RMSE gains on every dataset and baseline: prediction can remain competitive while group composition improves.

Component trade-offs are expected in a coupled correction/calibration system. Use \(=\) \texttt{NO\_CLEAR\_EFFECT} means no robust sealed E4 marginal effect, not that usability is conceptually irrelevant. Debt is statistically consistent but quantitatively small.

E2 Traffic supports the risk-semantic claim; E3 Traffic cross-vintage inference is \texttt{NOT\_COMPARABLE}. P2 warning windows are feasibility-audited; exact experimental optimality is not reconstructable. Formal real-mobility performance is not part of the current evidence.
